\chapter{Lógica matemática}
\section{Proposiciones}

\begin{lesson}[id=intro-proposiciones,weight=0]
Una proposición es un enunciado que puede ser verdadero o falso. Por ejemplo, \(2+2=4\) es verdadero, mientras que \(2+2=5\) es falso.
\end{lesson}

\begin{definition}[id=def-negacion,title={Negación},weight=0]
La negación de una proposición \(p\) se representa por \(\neg p\) y cambia su valor de verdad.
\end{definition}

\begin{activity}[id=respuesta-negacion,type=short,weight=1,due=2026-11-15]
Escribe la negación de la proposición: “Todos los números naturales son pares”.
\end{activity}

\begin{activity}[id=tabla-verdad,type=choice,weight=2]
¿Cuál es la negación de \(p\)?
\option{a}{\(p\)}
\option{b}{\(\neg p\)}
\option{c}{\(p\land q\)}
\correct{b}
\end{activity}

\begin{activity}[id=explicacion,type=long,weight=2]
Explica la diferencia entre una proposición y una pregunta cotidiana.
\end{activity}

\begin{deliverable}[id=entrega-logica,type=upload,weight=3,due=2026-12-01]
Entrega un PDF con una tabla de verdad y una explicación de sus resultados.
\end{deliverable}

\begin{table}[header=teal]
\row{Símbolo}{Lectura}
\row{\(\neg p\)}{No \(p\)}
\row{\(p\land q\)}{\(p\) y \(q\)}
\row{\(p\lor q\)}{\(p\) o \(q\)}
\end{table}
